\band{yellowhead}{Practice Midterm · Part 3 (ต่อ)}
\topic{4/5 Use Case: Grades ทั้งระบบ (L4, 64-65)}
\k{เริ่มวาด}: boundary+actors → goals จาก req → associations → include/extend ตามเงื่อนไข; \r{ไม่ใช่ flowchart}
\begin{center}\begin{tikzpicture}[x=1mm,y=.9mm,font=\fontsize{6.8}{7.1}\selectfont]
\draw[black!65](8,4)rectangle(38,-50);
\node[anchor=south,font=\fontsize{7.1}{7.4}\selectfont]at(23,4){Grade Submission};
\tikzset{uc/.style={draw,ellipse,align=center,inner sep=.8pt,text width=15mm,font=\fontsize{6.8}{7.1}\selectfont},outside/.style={align=center,text width=8mm,inner sep=.5pt,font=\fontsize{6.5}{6.8}\selectfont}}
\node[uc](login)at(23,0){Log in};
\node[uc](courses)at(23,-8){View courses\\and sections};
\node[uc](students)at(23,-16){Load registered\\students};
\node[uc](edit)at(23,-24){Edit / print\\grade sheet};
\node[uc](submit)at(23,-32){Submit grades};
\node[uc,text width=9mm](tr)at(15,-44){Update\\transcript};
\node[uc,text width=9mm](em)at(31,-44){Email\\students};
\draw(2.5,-14)circle(.7mm);\draw(2.5,-14.7)--(2.5,-17.5);\draw(1,-15.6)--(4,-15.6);\draw(2.5,-17.5)--(1,-19);\draw(2.5,-17.5)--(4,-19);
\node[outside,anchor=north]at(3,-20){Instructor};
\foreach\goal in{login,courses,students,edit,submit}{\draw(4,-15.6)--(\goal.west);}
\node[outside,draw](sso)at(43,0){\textless\textless actor\textgreater\textgreater\\KU SSO};
\node[outside,draw](reg)at(43,-12){\textless\textless actor\textgreater\textgreater\\KU\\Registration};
\node[outside,draw,text width=7.5mm](transcript)at(3.5,-44){\textless\textless actor\textgreater\textgreater\\KU\\Transcript};
\node[outside,draw](email)at(43,-44){\textless\textless actor\textgreater\textgreater\\KU Email};
\node[outside,draw,text width=16mm](info)at(23,-54){\textless\textless actor\textgreater\textgreater\\KU Student Info};
\draw(login.east)--(sso.west) (courses.east)--(reg.west) (students.east)--(reg.west) (tr.west)--(transcript.east) (em.east)--(email.west) (em.south)--(info.north);
\draw[dependency](submit.south west)--(tr.north);
\draw[dependency](submit.south east)--(em.north);
\node[fill=white,inner sep=.25pt,font=\fontsize{6.5}{6.8}\selectfont]at(15,-37){\textless\textless include\textgreater\textgreater};
\node[fill=white,inner sep=.25pt,font=\fontsize{6.5}{6.8}\selectfont]at(31,-37){\textless\textless include\textgreater\textgreater};
\end{tikzpicture}\end{center}
\extra ภาพจาก req; กล่อง actor = ระบบนอก; Student Info ให้ contacts; \b{Ready → submit → Submitted}; transcript/email บังคับ
\topic{5/5 State: object เดียว / event [guard] / effect}
\k{วิธีทำ}: object → states/events/[guards]/effects; initial/final/self-loop; \r{actor ไม่ใช่ state}; P = Programmer
\begin{center}\begin{tikzpicture}[x=1mm,y=1mm,font=\fontsize{6.8}{7.1}\selectfont]
\node[circle,fill,inner sep=.8pt](s)at(1,0){};
\node[state,text width=12mm](req)at(11,0){Requested};
\node[state,text width=12mm](approved)at(35,0){Approved};
\node[state,text width=12mm](work)at(35,-9){In Progress};
\node[state,text width=13mm](cc)at(11,-9){Closed by CCE};
\node[state,text width=19mm](cp)at(35,-18){Closed by Programmer};
\node[circle,draw,inner sep=1.1pt](e)at(11,-22){};\fill(11,-22)circle(.45mm);
\draw[flow](s)--node[above,font=\fontsize{6.2}{6.5}\selectfont]{save}(req);
\draw[flow](req)--node[above,align=center,font=\fontsize{6.2}{6.5}\selectfont]{approve\\{[CCE]}}(approved);
\draw[flow](req)--node[right,align=center,font=\fontsize{6.2}{6.5}\selectfont]{close\\{[CCE]}}(cc);
\draw[flow](approved)--node[left,font=\fontsize{6.2}{6.5}\selectfont]{start [P]}(work);
\draw[flow](work.east)--(46,-9)--(46,-5.5)--(38,-5.5)--([xshift=3mm]work.north);
\node[rotate=90,font=\fontsize{6.2}{6.5}\selectfont]at(47,-7){update [P]};
\draw[flow](work)--node[left,font=\fontsize{6.2}{6.5}\selectfont]{close [P]}(cp);
\draw[flow](cc)--(e);\draw[flow](cp.south)--(35,-22)--(e.east);
\end{tikzpicture}\end{center}
\k{Ticket}: save [System] → Requested; action อัปเดต Last Status/Due, แก้/ลบไม่ได้; closed ห้ามเพิ่ม; notes/due update เป็น self-loop
\k{เทียบ}: Registration Draft → Done เมื่อจ่าย / lock+email; เต็มยัง Draft. Grades Draft → Ready → Submitted; \r{ห้ามแต่ง state}
\topic{Git 20 คะแนน: เตรียมครบ Samples 2.1-2.9}
\k{โจทย์}: clean/คำสั่งสำเร็จ; conflict เฉพาะที่ระบุ. HEAD → branch → commit; origin/* = remote ที่รู้ล่าสุด; tag คงที่
{\smalltable\begin{tabularx}{\linewidth}{|C{4mm}|Y|}\hline
\rowcolor{lightrow}ข้อ & ฐาน → ผล / HEAD / ตัวชี้ที่ต้องรักษา\\\hline
2.1 & FF C→E; s=f=E, os=C; switch ท้าย → HEAD/main=B\\\hline
2.2 & HEAD=s=G, os=F; G ยังไม่ push; f=E แยกสาย → 3-way; graph ไม่ยืนยัน conflict\\\hline
2.3 & s:C→E (FF)→F; HEAD=s=F, f=E, os=C\\\hline
2.4 & s:D→E (FF)→F; HEAD=s=F, f=E, os=D\\\hline
2.5 & merge F/E→G (parents F/E); s=G, f=E, os=F; switch ท้าย → HEAD/main=B\\\hline
2.6 & main=B, origin/main=v1.0=H; switch ยัง B; FF→H; refs อื่นเดิม\\\hline
2.7 & HEAD=s=H (E/G), os=E; push f ที่ G ไม่ส่ง H; push s → os=H\\\hline
2.8 & s:C→E; สร้าง/switch f ใหม่→F; s=feature/pass\_rate=E, os=C\\\hline
2.9 & s:D→G (FF)→H (G/F); push→os=H; HEAD=s=H, f=F\\\hline
\end{tabularx}}
\k{ชื่อย่อ}: s=staging, f=feature, os=origin/staging; main/origin/main=B ยกเว้น FF ใน 2.6; ใช้ชื่อจริงตามข้อ
\topic{Ancestry ย่อเฉพาะช่วงเปลี่ยน; refs ดูตาราง}
\begin{center}\begin{tikzpicture}[x=1mm,y=1mm,font=\fontsize{7.1}{7.4}\selectfont]
\foreach\lab/\x in{C/2,D/8,E/14,F/20}{\node[circle,draw,inner sep=1pt](left\lab)at(\x,0){\lab};}
\draw(leftC)--(leftD)--(leftE)--(leftF);
\node[anchor=south,inner sep=.5pt]at(11,2){2.3 FF + commit F};
\node[anchor=north,inner sep=.5pt]at(11,-2){staging=F, feature=E};
\node[circle,draw,inner sep=1pt](d)at(27,0){D};\node[circle,draw,inner sep=1pt](f)at(35,0){F};\node[circle,draw,inner sep=1pt,fill=minthead!45](g)at(45,0){G};\node[circle,draw,inner sep=1pt](e)at(35,-5){E};
\draw(d)--(f)--(g);\draw(d)--(e)--(g);
\node[anchor=south,inner sep=.5pt]at(36,2){2.5 3-way};
\node[anchor=north,inner sep=.5pt]at(36,-7){G parents F/E};
\end{tikzpicture}\end{center}
\topic{คำสั่ง: b/ref/files/msg ใช้ค่าจากโจทย์}
\par\code{git status}: ดู branch/staged/unstaged/untracked; ไม่แก้
\par\code{git diff}: เทียบ working tree กับ index เฉพาะ tracked; ไม่แสดงไฟล์ untracked
\par\code{git diff --cached}: ดู staged: index เทียบ HEAD
\par\code{git add <files>}: เก็บเวอร์ชันไฟล์ที่ระบุเข้า index; ยังไม่ commit หรือเลื่อน branch
\par\code{git commit -m "<msg>"}: บันทึก index เป็นลูก HEAD; เลื่อน current branch/HEAD; ไม่ push
\par\code{git branch <b>}: สร้าง b ที่ HEAD; HEAD/ไฟล์เดิม
\par\code{git switch <b>}: HEAD/ไฟล์ไป b เมื่อ clean; ตัวชี้ branch ไม่เลื่อน
\par\code{git merge <ref>}: รวมเข้า current; FF เลื่อน ref ไม่สร้าง commit; แยกสายสร้าง 2 parents; source เดิม
\par\code{git merge --ff-only origin/<b>}: เลื่อน current แบบ FF เท่านั้น; แยกสายแล้วปฏิเสธ
\par\code{git fetch origin}: รับ commits/อัปเดต origin/*; ไม่เลื่อน local branch/HEAD หรือไฟล์
\par\code{git push origin <b>}: ส่ง b ไป remote/อัปเดต origin/b เมื่อสำเร็จ; HEAD/ไฟล์เดิม
\k{2.3}: staging = sprint\_1\_staging, feature/load\_scores; แก้ score\_data.py, scores.csv; สร้าง validation\_notes.txt; message \code{"Document score validation cases"}
\k{2.8}: s=sprint\_2\_staging; merge feature/pass\_rate; branch+switch feature/grade\_distribution; แก้ report.py, scores.csv; เพิ่ม grade\_ranges.csv; message \code{"Add grade distribution report"}
\k{2.9}: b=sprint\_2\_staging, ref=feature/export\_results; switch b → FF origin/b → merge ref → push b
\r{ห้าม git add .}; วาด parent/local/origin/tag/HEAD ทุกขั้น
